
Experiments test whether the two components of bidirectional alignment---user learning and AI responsiveness---exist in RLHF-trained conversational systems. These are existence tests, not comparative evaluations against baselines.

\subsubsection{Dataset}

The HH-RLHF dataset \citep{bai2022constitutional} contains 161,000 multi-turn conversations between humans and an RLHF-trained AI assistant. Each conversation is rated by human annotators on helpfulness and harmlessness dimensions (scale 0-1). Conversations vary in length (2-20+ turns), topic (open-domain), and structure.

HH-RLHF is well-suited for studying bidirectional alignment because: (1) it contains genuine multi-turn interactions where users can learn AI capabilities over successive queries, (2) the AI system was trained with RLHF and thus may exhibit responsiveness to user patterns, and (3) helpfulness ratings allow future correlation analysis between alignment metrics and conversation quality.

\textbf{Filtering criteria}:
\begin{enumerate}
\item Minimum turn count: Conversations with $\geq 5$ turns. Reformulation slope estimation requires at least 4-5 query pairs for linear regression reliability.
\item Reformulation coverage: For Hypothesis 1, further filtering to conversations containing at least one detected reformulation. Conversations where all queries are novel yield null slopes and are excluded.
\end{enumerate}

After filtering, 88 conversations were analyzed for Hypothesis 1 and 169,352 conversations for Hypothesis 2. The large sample size for Hypothesis 2 reflects lenient filtering requirements, while the small sample for Hypothesis 1 reflects strict reformulation detection thresholds.

\subsubsection{Experimental Questions}

\textbf{Hypothesis 1 (h-e1): User Learning Exists}  
\textit{Research question}: Do users reformulate queries less frequently as conversations progress?  
\textit{Operationalization}: Mean reformulation slope across conversations is significantly negative ($\text{mean slope} < 0, p < 0.05$).  
\textit{Prediction}: If users learn which queries the AI handles effectively, they reformulate less over turns. Negative slope provides evidence of learning curve effects.

\textbf{Hypothesis 2 (h-e2): AI Responsiveness Exists}  
\textit{Research question}: Does AI response diversity correlate with user query diversity?  
\textit{Operationalization}: Pearson correlation between query diversity and response diversity is significantly positive ($r > 0.35, p < 0.05$).  
\textit{Prediction}: If the RLHF-trained AI exhibits behavioral responsiveness, diverse queries should elicit diverse responses. Positive correlation provides evidence of responsiveness.

These are conservative tests. User learning predicting task success is not claimed (that requires regression analysis, planned for future work). A threshold of $r > 0.35$ rather than the initially proposed $r > 0.4$ is set, acknowledging that lexical diversity may underestimate semantic responsiveness.

\subsubsection{Experimental Procedure}

\textbf{Hypothesis 1: Reformulation Slope Detection}

For each conversation $c$ with $T \geq 5$ turns:
\begin{enumerate}
\item Extract query pairs: For turns $t = 2, \ldots, T$, extract consecutive query pairs $(q_{t-1}, q_t)$.
\item Compute reformulation indicators: For each pair, compute semantic similarity ($\text{SBERT}\_\text{cosine}$ using all-MiniLM-L6-v2 model), syntactic distance (normalized Levenshtein distance), and reformulation label ($r_t = 1$ if $\text{sim} > 0.7$ AND $\text{dist} > 0.3$).
\item Compute slope: Fit linear regression $r_t \sim \beta_0 + \beta_1 \cdot t$. Extract slope coefficient $\beta_1$.
\item Aggregate: Collect slopes from all conversations. Test $H_0$: $\text{mean}(\beta_1) = 0$ vs $H_1$: $\text{mean}(\beta_1) < 0$ using one-sample t-test ($\alpha = 0.05$).
\end{enumerate}

\textbf{Hypothesis 2: Diversity Correlation}

For each conversation $c$:
\begin{enumerate}
\item Compute query diversity: $D_{\text{query}}(c) = \frac{|\text{unique unigrams in all queries}|}{|\text{total unigrams in all queries}|}$.
\item Compute response diversity: $D_{\text{response}}(c) = \frac{|\text{unique unigrams in all responses}|}{|\text{total unigrams in all responses}|}$.
\item Aggregate: Collect $(D_{\text{query}}, D_{\text{response}})$ pairs from all conversations. Compute Pearson correlation $r$ with significance test ($\alpha = 0.05$).
\end{enumerate}

\subsubsection{Baselines}

No baselines are compared in this paper. Experiments are existence tests (do these signals occur?) rather than comparative evaluations. Future work will test whether coupling strength predicts helpfulness beyond individual metrics, requiring baseline comparisons against unidirectional metrics and alternative coupling formulations.

\subsubsection{Evaluation Metrics}

\textbf{Primary Metrics}:
\begin{itemize}
\item Hypothesis 1: Mean reformulation slope, p-value (one-sample t-test), Cohen's d effect size.
\item Hypothesis 2: Pearson correlation $r$, p-value (correlation test), 95\% confidence interval.
\end{itemize}

\textbf{Success Criteria}:
\begin{itemize}
\item Hypothesis 1: Mean slope significantly negative ($p < 0.05$) with $|d| > 0.2$ (small-to-medium effect).
\item Hypothesis 2: Correlation $r > 0.35$ with $p < 0.05$ and confidence interval excluding zero.
\end{itemize}

\subsubsection{Reproducibility}

All experiments use fixed random seed (42) for dataset sampling and SBERT embedding initialization. HH-RLHF is publicly available at https://huggingface.co/datasets/Anthropic/hh-rlhf. Reformulation detection thresholds ($\text{sim} > 0.7, \text{dist} > 0.3$) were preregistered in the verification plan.
